\section{Modul 01. Keluaran}

\subsection{Tujuan}

Setelah menyelesaikan modul ini, mahasiswa mampu:
\begin{enumerate}[nosep]
    \item mencetak satu karakter dengan \texttt{INT 21h} fungsi
          \texttt{AH=02h};
    \item mencetak satu kalimat dengan \texttt{INT 21h} fungsi
          \texttt{AH=09h};
    \item menyusun string dengan direktif \texttt{DB} yang diakhiri
          \texttt{\$} dan memuat \texttt{0Dh}, \texttt{0Ah};
    \item menjelaskan perbedaan penyiapan \texttt{DS} pada program
          \texttt{.COM} dan \texttt{.EXE};
    \item merakit, menautkan, dan menjalankan kedua program dengan
          TASM dan TLINK.
\end{enumerate}

\subsection{Dasar singkat}

\subsubsection*{Layanan 02h: mencetak satu karakter}

Layanan \texttt{AH=02h} mencetak satu karakter ke layar. Sebelum
\texttt{INT 21h}, nomor layanan diletakkan pada \texttt{AH} dan
karakter yang dicetak diletakkan pada \texttt{DL}. Karena
\texttt{DL} hanya 8 bit, satu panggilan layanan ini hanya membawa
satu karakter. Karakter dapat ditulis sebagai literal, misalnya
\texttt{'A'}, atau sebagai kode ASCII, misalnya \texttt{41h}.
Setelah karakter tercetak, kursor bergerak satu posisi ke kanan.
Layanan ini tidak mengubah warna dan tidak memakai perulangan.

\subsubsection*{Layanan 09h: mencetak satu kalimat}

Layanan \texttt{AH=09h} mencetak deretan karakter yang ditunjuk
pasangan \texttt{DS:DX}. Sebelum \texttt{INT 21h}, \texttt{DX} diisi
alamat offset string dan \texttt{AH} diisi \texttt{09h}. String
didefinisikan pada segmen data dengan tipe \texttt{DB}, lalu diakhiri
tanda \texttt{\$}. Tanda \texttt{\$} adalah penanda akhir dan tidak
ikut tercetak; DOS berhenti mencetak begitu menemukannya.

Karakter \texttt{0Dh} mengembalikan kursor ke awal baris, sedangkan
\texttt{0Ah} menurunkannya satu baris. Bila keduanya diletakkan
sebelum \texttt{\$}, kalimat berikutnya dimulai pada baris baru.
Ringkasan karakter yang dipakai pada modul ini ditunjukkan pada
Tabel~\ref{tab:m01-kontrol}.

\begin{table}[h]
    \centering
    \caption{Karakter kendali dan penanda akhir string.}
    \label{tab:m01-kontrol}
    \small
    \begin{tabular}{@{}lll@{}}
        \toprule
        \textbf{Nilai} & \textbf{Nama} & \textbf{Akibat} \\
        \midrule
        \texttt{0Dh} & \emph{carriage return} &
        kursor kembali ke awal baris \\
        \texttt{0Ah} & \emph{line feed} &
        kursor turun satu baris \\
        \texttt{24h} & tanda \texttt{\$} &
        akhir string; tidak ikut dicetak \\
        \bottomrule
    \end{tabular}
\end{table}

Satu string boleh ditulis dalam beberapa baris \texttt{DB}; assembler
menyimpannya berurutan. Tanda \texttt{\$} cukup ditulis sekali, pada
potongan terakhir.

\subsubsection*{Register, segmen, dan pemilihan layanan}

Layanan \texttt{02h} dan \texttt{09h} sama-sama dipanggil dengan
\texttt{INT 21h}, tetapi register yang disiapkan berbeda. Perbedaan
itu dirangkum pada Tabel~\ref{tab:m01-layanan}. Nomor layanan pada
\texttt{AH} harus ditulis sebelum setiap panggilan.

\begin{table}[h]
    \centering
    \caption{Persiapan register untuk layanan keluaran.}
    \label{tab:m01-layanan}
    \small
    \begin{tabular}{@{}lll@{}}
        \toprule
        \textbf{Layanan} & \textbf{Register masukan} &
        \textbf{Keluaran} \\
        \midrule
        \texttt{AH=02h} & \texttt{DL} = satu karakter &
        satu karakter di layar \\
        \texttt{AH=09h} & \texttt{DS:DX} = alamat string &
        kalimat sampai tanda \texttt{\$} \\
        \texttt{AH=4Ch} & \texttt{AL} = kode keluar &
        program berakhir \\
        \bottomrule
    \end{tabular}
\end{table}

Pada program \texttt{.COM}, DOS sudah menetapkan \texttt{DS} ke
segmen program. Pada program \texttt{.EXE}, \texttt{DS} belum
menunjuk segmen data, sehingga harus diisi dengan
\texttt{MOV AX,@DATA} lalu \texttt{MOV DS,AX}. Operator
\texttt{OFFSET} menghasilkan alamat offset variabel untuk \texttt{DX}.

Bila tanda \texttt{\$} lupa ditulis, DOS terus mencetak isi memori di
belakang string. Pastikan \texttt{\$} berada sesudah \texttt{0Dh},
\texttt{0Ah} yang ingin dicetak.

Rujukan singkat: Gunawan, Bagian 6.1 dan 6.4; S'to, Bab VII dan XII;
buku ajar, berkas \texttt{hello.asm} dan \texttt{chapters/07}.

\subsection{Alat dan bahan}

\begin{itemize}[nosep]
    \item Komputer dengan emu8086 terpasang.
    \item DOSBox dengan TASM, TLINK, dan DEBUG.
    \item Berkas \texttt{code/m01-cetak-karakter.asm} dan
          \texttt{code/m01-cetak-kalimat.asm}.
    \item Lembar jurnal praktikum.
\end{itemize}

\subsection{Contoh program}

Program pertama mencetak satu karakter dengan layanan \texttt{02h}.
Program ini berbentuk \texttt{.COM}, sehingga tidak mengisi
\texttt{DS} dan tidak memerlukan pola \texttt{JMP MULAI}.

\lstinputlisting[
    caption={Mencetak satu karakter dengan \texttt{AH=02h}.},
    label={lst:m01-karakter}
]{code/m01-cetak-karakter.asm}

Program kedua mencetak dua baris dengan layanan \texttt{09h}.
Program ini berbentuk \texttt{.EXE}, sehingga \texttt{DS} diisi lebih
dahulu; \texttt{PESAN} memuat \texttt{0Dh}, \texttt{0Ah} sebelum
tanda \texttt{\$}.

\lstinputlisting[
    caption={Mencetak kalimat dengan \texttt{AH=09h}.},
    label={lst:m01-kalimat}
]{code/m01-cetak-kalimat.asm}

\subsection{Langkah kerja}

\begin{enumerate}
    \item Siapkan folder kerja, lalu salin kedua berkas
          \texttt{.asm} ke dalamnya.
    \item Rakit dan tautkan program pertama:
          \texttt{tasm m01-cetak-karakter.asm}, lalu
          \texttt{tlink /t m01-cetak-karakter.obj}. Jalankan
          hasilnya dan catat keluarannya.
    \item Rakit dan tautkan program kedua. Jalankan kedua perintah
          berikut pada baris terpisah:
          \begin{quote}
          \texttt{tasm m01-cetak-kalimat.asm}\\
          \texttt{tlink m01-cetak-kalimat.obj}
          \end{quote}
          Jalankan hasilnya dan catat keluarannya.
    \item Ubah karakter pada program pertama dari \texttt{'A'} menjadi
          huruf lain, rakit ulang, lalu jalankan. Jelaskan bagian mana
          yang berubah.
    \item Ubah salah satu kalimat pada \texttt{PESAN}, lalu tambahkan
          satu pasang \texttt{0Dh}, \texttt{0Ah} sebelum \texttt{\$}.
          Rakit ulang dan bandingkan keluarannya.
    \item Periksa program kedua dengan DEBUG: jalankan
          \texttt{debug m01-cetak-kalimat.exe}, ketik \texttt{r},
          lalu \texttt{d ds:0 L 40}. Temukan string \texttt{PESAN}
          dan byte \texttt{24h} pada hasil pemeriksaan, lalu keluar
          dengan \texttt{q}.
\end{enumerate}

\subsection{Tugas}

\begin{enumerate}
    \item Jelaskan tugas \texttt{AH} dan \texttt{DL} pada layanan
          \texttt{02h}.
    \item Jelaskan mengapa layanan \texttt{09h} memerlukan
          \texttt{DS:DX}, sedangkan \texttt{02h} tidak.
    \item Tentukan keluaran yang terjadi bila tanda \texttt{\$}
          dihapus dari string, lalu jelaskan alasannya.
    \item Jelaskan perbedaan penyiapan \texttt{DS} pada
          \texttt{.COM} dan \texttt{.EXE}.
    \item Ubah program kedua agar mencetak nama lengkap Anda pada
          baris pertama dan nama modul pada baris kedua. Tentukan
          letak tanda \texttt{\$} yang benar.
\end{enumerate}

\subsection{Hasil yang dikumpulkan}

\begin{enumerate}
    \item Kedua berkas \texttt{.asm} beserta versi modifikasinya.
    \item Tangkapan layar keluaran program pertama dan program kedua.
    \item Tangkapan layar DEBUG yang memperlihatkan string dan byte
          \texttt{24h}.
    \item Lembar jurnal berisi jawaban langkah kerja dan tugas.
\end{enumerate}
